%%%%%%%%%%%%%%%%%%%%%%%%%%%%%%%%%%%%%%%%%%%%%%%%%%%%%%%%%%%%%%%%%%%%%%%%%%%%%%%%%%%%%%%%%%%%%%%%%%%%%%%%%%%%%%%%
% Sablona Bc+Mgr+RNDr (CZ) pro PriF MU                                                                       %%%
% Autor: Petr Zemanek (zemanekp@math.muni.cz)                                                                %%%
% Licence: tento soubor je k dispozici bez jakychkoli omezeni                                                %%%
% Pripominky, dotazy, namety smerujte na diskusni forum: https://is.muni.cz/auth/discussion/sablona-prif/    %%%
% Typeset in LaTeX-2e                                                                                        %%%
% Verze: 2.2 (11. prosince 2019)                                                                             %%%
%%%%%%%%%%%%%%%%%%%%%%%%%%%%%%%%%%%%%%%%%%%%%%%%%%%%%%%%%%%%%%%%%%%%%%%%%%%%%%%%%%%%%%%%%%%%%%%%%%%%%%%%%%%%%%%%

% \documentclass[11pt,a4paper,oneside,final]{book} %% PRO JEDNOSTRANNY TISK
% \documentclass[11pt,a4paper,twoside,final]{book} %% PRO OBOUSTRANNY TISK
% \documentclass[12pt,a4paper,oneside,final]{book} %% PRO JEDNOSTRANNY TISK
\documentclass[12pt,a4paper,twoside,final]{book} %% PRO OBOUSTRANNY TISK

%%%%%%%%%%%%%%%%%%%%%%%%%%%%%%%%%%%%%%%%%%%%%%%%%%%%%%%%%%%%%%%%%%%%%%%%%%%%%%%%%%%%%%%%%%%%%%%%%%%%%%
%%%%%%%%%%%%%%%%%%%%%%%%%%%%%%%%%%%%%% ZAKLADNI NASTAVENI %%%%%%%%%%%%%%%%%%%%%%%%%%%%%%%%%%%%%%%%%%%%

%%%%%%%%%%%%
%% Zvolte kodovani dokumentu
%%%%%%%%%%%%

\usepackage[utf8]{inputenc} 
% \usepackage[cp1250]{inputenc} %% NASTAVENI PRO WINDOWS
% \usepackage[T1]{fontenc} %% puvodni evropske fonty->spatne umistene hacky nad c, d s hackem ma za sebou mezeru
\usepackage[IL2]{fontenc} %fonty vyladene pro cestinu/slovenstinu

%%%%%%%%%%%%
%% Nastavte si jazyk dokumentu 
%% lze pouzit volbu babel (pro pdflatex a latex) nebo bez babelu (pro pdfcslatex a cslatex
%%
%% pri volbe bez babelu bude zlobit v ukazkovych souborech prikaz ``\shorthandoff{-}'' a ``\shorthandoff{-}''
%% ty staci vymazat a je to bez problemu
%%
%% pri volbe ``\usepackage[english]{babel}'' nefunguji ceske uvozovky pomoci \uv{..} - staci nahradit ``..''
%%%%%%%%%%%%


\usepackage[english, czech]{babel}
% \usepackage{czech}
% \usepackage{slovak}

%%%%%%%%%%%%
%% Muzete si zkusit zmenit font dokumentu - pak ale musite v souboru zmenit nektera nastaveni (nebot je nastaveno
%% \overfullrule10pt je preteceni textu ihned vyznaceno cernym obdelnikem)
%% Veskere nastaveni je testovano pro pouzitou volbu
%%%%%%%%%%%%

\usepackage{fourier}
% \usepackage{mathpazo}
% \usepackage{newcent}
% \usepackage{palatino}
% \usepackage{utopia}
% \usepackage{mathptmx}  %% volne dostupny font Adobe Times Roman
%  \usepackage[mtbold,mtplusscr,mtpluscal]{mathtime} % komercni matematicky font, dostupne pouze na UMS

%%%%%%%%%%%%%%%%%%%%%%%%%%%%%%%%%%%%%%%%%%%%%%%%%%%%%%%%%%%%%%%%%%%%%%%%%%%%%%%%%%%%%%%%%%%%%%%%%%%%
%%%%%%%%%%%%%%%%%%%%%%%%%%%%%%%% BALIKY POTREBNE PRO SABLONU %%%%%%%%%%%%%%%%%%%%%%%%%%%%%%%%%%%%%%%

\usepackage{longtable, lipsum} 
%%% balik ``longtable`` je potrebny pro sazeni kapitoly s pouzitym znacenim, jinak neni potreba
\usepackage{xcolor} %% definice muni modre barvy

%%%%%%%%%%%%%%%%%%%%%%%%%%%%%%%%%%%%%%%%%%%%%%%%%%%%%%%%%%%%%%%%%%%%%%%%%%%%%%%%%%%%%%%%%%%%%%%%%%%%
%%%%%%%%%%%%%%%%%%%%%%%%%%%%% BALIKY POTREBNE PRO MATEMATIKU %%%%%%%%%%%%%%%%%%%%%%%%%%%%%%%%%%%%%%%

\usepackage{amsmath,amssymb,amsthm}

%%%%%%%%%%%%%%%%%%%%%%%%%%%%%%%%%%%%%%%%%%%%%%%%%%%%%%%%%%%%%%%%%%%%%%%%%%%%%%%%%%%%%%%%%%%%%%%%%%%%
%%%%%%%%%%%%%%%%%%%%%%%%%%%%% BALIKY POTREBNE PRO MNE %%%%%%%%%%%%%%%%%%%%%%%%%%%%%%%%%%%%%%%

\usepackage[backend=biber,citestyle=authoryear-comp,natbib=true,sorting=nyt,maxcitenames=2,doi=true]{biblatex}
\let\cite\citep
\addbibresource{../../DP_VamberskyT.bib}
\usepackage{csquotes}
\AtBeginBibliography{\selectlanguage{english}}
\AtEveryBibitem{%
  \clearfield{issue}%
}

%%%%%%%%%%%%%%%%%%%%%%%%%%%%%%%%%%%%%%%%%%%%%%%%%%%%%%%%%%%%%%%%%%%%%%%%%%%%%%%%%%%%%%%%%%%%%%%%%
%%%%%%%%%%%%%%%%%% NASTAVENI POVINNYCH CASTI DLE SMERNICE DEKANA %%%%%%%%%%%%%%%%%%%%%%%%%%%%%%%%

%%%%%%%%%%%%%%%%%%%%%%%%%%%%%%%%%%%%%%%%%%%%%%%%%%%%%%
%% NASTAVTE KONKRETNI UDAJE V CESTINE A ANGLICTINE  %%
%%%%%%%%%%%%%%%%%%%%%%%%%%%%%%%%%%%%%%%%%%%%%%%%%%%%%%

\usepackage[Mgr,Barevne]{sci.muni.thesis}
%% Mozne volby:
%% Bc - pro bakalarskou praci
%% Mgr - pro diplomovou praci
%% RNDr - pro rigorozni praci
%% Barevne - pro dokument s barevnymi odkazy
%% Tisk - pro dokument v cernobile barve

\NazevUstavu{Ústav biochemie}{Department of Biochemistry}

\RokOdevzdaniPrace{2026}

\AkademickyRok{2025/2026}

\Autor{Tomáš Vamberský}{Bc.\,Tomáš Vamberský}

\NazevPrace{Testování efektu AI Foundational models na analýzu NGS dat DNA sekvenování}{Testování efektu AI Foundational models na analýzu NGS dat DNA sekvenování}{Testing AI Foundational Models for Impact on NGS DNA Sequencing Data Analysis}


\VedouciPraceSTituly{Mgr.\,Vojtěch Bystrý, Ph.D.}

\StudijniProgram{Biochemie}{Biochemistry}

\StudijniObor{Bioinformatika}{Bioinformatics}

\PocetStran{??\,$+$\,??}

\KlicovaSlova{AI; Foundation model; Transcription factor}{AI; Foundation model; Transkripční faktor}

\Abstrakty%
{V~této bakalářské/diplomové/rigorózní práci se věnujeme ...}%
{In this thesis we study ...}

\TextPodekovani%
{Na tomto místě bych chtěl(-a) poděkovat ...}

\TextProhlaseni%
{Prohlašuji, že jsem svouji diplomovou práci vypracoval samostatně pod vedením vedoucího práce 
s~využitím informačních zdrojů, které jsou v~práci citovány. Prohlašuji, že jsem pro stylistické úpravy textu využil nástroje AI a to plně v~souladu s~principy akademické integrity.}

\DatumProhlaseni{13. Května 2026}



\newcommand{\Cbb}{\mathbb{C}}
\newcommand{\Rbb}{\mathbb{R}}
\newcommand{\Zbb}{\mathbb{Z}}
\newcommand{\Nbb}{\mathbb{N}}
\DeclareMathOperator{\tg}{tg}
\newcommand{\mybeta}{$\beta$}
\newcommand{\mygamma}{$\gamma$}
\newcommand{\myepsilon}{$\varepsilon$}
\newcommand{\myzeta}{$\zeta$}
\newcommand{\myeta}{$\eta$}
\newcommand{\mytheta}{$\theta$}
\newcommand{\mysigma}{$\sigma$}
\newcommand{\fix}[1]{\textit{#1}}
%%%%%%%%%%%%%%%%%%%%%%%%%%%%%%%%%%%%%%%%%%%%%%%%%%%%%%%%%%%%%%%%%%%%%%%%%%%%%%%%%%%%%%%%%%%%%%%%%%%%%%
%%%%%%%%%%%%%%%%%%%%%%%% VYTVOR POMOCNY SOUBOR PRO REJSTRIK %%%%%%%%%%%%%%%%%%%%%%%%%%%%%%%%%%%%%%%%%%

\makeindex

%%%%%%%%%%%%%%%%%%%%%%%%%%%%%%%%%%%%%%%%%%%%%%%%%%%%%%%%%%%%%%%%%%%%%%%%%%%%%%%%%%%%%%%%%%%%%%%%%%%%%%
%%%%%%%%%%%%%%%%%%%%%%% ZAPNUTI OPAKOVANI MATEMATICKYCH SYMBOLU %%%%%%%%%%%%%%%%%%%%%%%%%%%%%%%%%%%%%%

% \input{opakuj.sty}

%%%%%%%%%%%%%%%%%%%%%%%%%%%%%%%%%%%%%%%%%%%%%%%%%%%%%%%%%%%%%%%%%%%%%%%%%%%%%%%%%%%%%%%%%%%%%%%%%%%%%%
%%%%%%%%%%%%%%%%%%%%%%%%%%%%%%%%%% ZACATEK DOKUMETU %%%%%%%%%%%%%%%%%%%%%%%%%%%%%%%%%%%%%%%%%%%%%%%%%%
%%%%%%%%%%%%%%%%%%%%%%%%%%%%%%%%%%%%%%%%%%%%%%%%%%%%%%%%%%%%%%%%%%%%%%%%%%%%%%%%%%%%%%%%%%%%%%%%%%%%%%

\begin{document}

%%%%%%%%%%%%%%%%%%%%%%%%%%%%%%%%%%%%%%%%%%%%%%%%%%%%%%%%%%%%%%%%%%%%%%%%%%%
%%%%%%%%%%%% POVINNE CASTI PRO Bc. A Mgr. PRACE V CESTINE %%%%%%%%%%%%%%%%%

\VytvorPovinneStrany

%%%%%%%%%%%%%%%%%%%%%%%%%%%%%%%%%%%%%%%%%%%%%%%%%%%%%%%%%%%%%%%%%%%%%%%%%%%%%%%%%%%%%%%%%%%%%%%%%%%
%%%%%%%%%%%%%%%%%%%%%%%%%%%%%%%%%% ABSTRAKT + ABSTRACT %%%%%%%%%%%%%%%%%%%%%%%%%%%%%%%%%%%%%%%%%%%%

%\AbstraktyNaJedneStrane

%%%%%%%%%%%%%%%%%%%%%% POKUD SE ABSTRAKTY NEVEJDOU NA JEDNU STRANU, POUZIJTE TOTO NASTAVENI %%%%%%%

\AbstraktyNaDvouStranach


%%%%%%%%%%%%%%%%%%%%%%%%%%%%%%%%%%%%%%%%%%%%%%%%%%%%%%%%%%%%%%%%%%%%%%%%%%%%%%%%%%%%%%%%%%%%%%%%%%%
%%%%%%%%%%%%%%%%%%%%%%%%%%%%%%%%%%%%%% SEM VLOZIT ZADANI %%%%%%%%%%%%%%%%%%%%%%%%%%%%%%%%%%%%%%%%%%
%%% OSKENUJTE JEJ A VE FORMATU PDF VLOZTE DO STEJNE SLOZKY (SOUBOR POJMENUJTE NAPR. zadani.pdf) %%%
%%%%%% PAK ''ZAPROCENTUJTE'' RADEK S \SemVlozitZadani A ODPROCENTUJTE RADEK S \VlozZadani %%%%%%%%%
%%%%%%%%%%%%%%%%%%%%%%%%%%%%%%%%%%%%%%%%%%%%%%%%%%%%%%%%%%%%%%%%%%%%%%%%%%%%%%%%%%%%%%%%%%%%%%%%%%%

\SemVlozitZadani

% \VlozZadani{NazevSouboruSeZadanim}
% \VlozZadani{zadani.pdf}

%%%%%%%%%%%%%%%%%%%%%%%%%%%%%%%%%%%%%%%%%%%%%%%%%%
%%%%%%%%%% PODEKOVANI + PROHLASENI %%%%%%%%%%%%%%%

\PodekovaniAProhlaseni
%
% \ProhlaseniBezPodekovani
%
%%%%%%%%%%%%%%%%%%%%%%%%%%%%%%%%%%%%%%%%%%%%%%%%%%
%%%%%%%%%%%%%%%%%%% OBSAH %%%%%%%%%%%%%%%%%%%%%%%%

\pdfbookmark{Contents}{Contents}

\renewcommand{\contentsname}{Table of Contents}
\renewcommand{\listfigurename}{List of Figures}
\renewcommand{\listtablename}{List of Tables}

\VytvorObsah
\clearpage

\addcontentsline{toc}{chapter}{\listfigurename}
\listoffigures
\clearpage

\addcontentsline{toc}{chapter}{\listtablename}
\listoftables
\clearpage

%%%%%%%%%%%%%%%%%%%%%%%%%%%%%%%%%%%%%%%%%%%%%%%%%%%%%%%%%%%%%%%%%%%%%%%%%%%%%%%%%%%%%%%%%%%%%%%%%%%
%%%%%%%%%%%%%%%%%%%%%%%%%%%%%%%%%%% TEXT PRACE %%%%%%%%%%%%%%%%%%%%%%%%%%%%%%%%%%%%%%%%%%%%%%%%%%%%
\renewcommand{\chaptername}{Chapter}

\HlavickaZnaceni
\vloz{text_prace/01_Znaceni}
\cleardoublepage

\HlavickaUvod
\setcounter{page}{1}
\pagenumbering{arabic}
\vloz{text_prace/02_Uvod}
\cleardoublepage

\renewcommand{\chaptermark}[1]{\markboth{\thechapter. #1}{}}
\renewcommand{\sectionmark}[1]{\markright{\thesection. #1}{}}
\HlavickaKapitoly
\vloz{text_prace/03_Teorie}
\cleardoublepage

\vloz{text_prace/04_Cile}

\vloz{text_prace/05_Metody}
\cleardoublepage

\vloz{text_prace/06_Vysledky_Diskuse}
\cleardoublepage

\HlavickaZaver
\vloz{text_prace/07_Zaver}
\cleardoublepage


\renewcommand{\bibname}{Bibliography}
%%\nocite{*} % Toto automaticky vygeneruje seznam vsech citaci ze souboru Zdroje.bib. Smaz pokud chces pouze seznam citaci, ktere se objevuji v textu
\printbibliography
\addcontentsline{toc}{chapter}{\bibname}
\cleardoublepage

\HlavickaPriloha 
\vloz{text_prace/09_Priloha}
%%
%% pro vytvoreni rejstriku se spravny ceskym razenim pouzijte 
%% csindex -d -h -k -z il2 nazev_souboru.idx
%%
%% nebo 
%% texindy -I latex -L czech -M lang/czech/utf8 nazev_souboru.idx
%%
%% na Ustavu matematiky a statistiky zadejte
%% /opt/texlive/2010/bin/i386-linux/texindy -I latex -L czech -M lang/czech/utf8 nazev_souboru.idx

%%%%%%%%%%%%%%%%%%%%%%%%%%%%%%%%%%%%%%%%%%%%%%%%
%%%%%%%%%%% PRAZDNA STRANA NA ZAVER %%%%%%%%%%%%
%%%%%%%%%%%%%%%%%%%%%%%%%%%%%%%%%%%%%%%%%%%%%%%%

\newpage
\thispagestyle{empty}
\fancyhf{}
\newpage
\mbox{}

\end{document}
